% Teil II: der durchgerechnete Weg.

% Gedankenschritt: negative Exponenten sind Kehrwerte, gebrochene Wurzeln;
% die Potenzgesetze reduzieren Mehrfach-Potenzen auf eine Potenz.

\subsection*{Gegeben und gesucht}
(a)--(d) Werte ohne Taschenrechner; (e)--(h) Umformen bzw. vereinfachen.

\subsection*{Ansatz}
\textbf{Alle Teile beruhen auf zwei Kernideen:} (1) ein negativer
Exponent kehrt den Kehrwert um, (2) ein gebrochener Exponent ist eine
Wurzel. Wo mehrere Potenzen gemischt sind, wendet man die
Potenzgesetze rückwärts an, um auf eine einzelne Potenz zu kommen.

\subsection*{a) $2^{-3}$}
\begin{equation}
  2^{-3}=\frac1{2^3}=\frac18 .
\end{equation}
Gedanke: nicht ``minus drei Zweien``, sondern ``Kehrwert von
$2^3=8$``.

\subsection*{b) $16^{1/2}$}
\begin{equation}
  16^{1/2}=\sqrt{16}=4 .
\end{equation}

\subsection*{c) $27^{-1/3}$}
\begin{equation}
  27^{-1/3}=\frac1{27^{1/3}}=\frac1{\sqrt[3]{27}}=\frac13 .
\end{equation}

\subsection*{d) $\bigl(\frac49\bigr)^{-1/2}$}
\begin{equation}
  \Bigl(\frac49\Bigr)^{-1/2}=\Bigl(\frac94\Bigr)^{1/2}=\frac32 .
\end{equation}
\textbf{Gedankenschritt:} Zwei Umkehr-Operationen hintereinander: der
negative Exponent kehrt den Bruch um ($\tfrac49\to\tfrac94$), dann zieht
der $\tfrac12$-Exponent die Wurzel.

\subsection*{e) $\frac1{a^{-5}}$ als Potenz mit positivem Exponenten}
\begin{equation}
  \frac1{a^{-5}}=a^5 .
\end{equation}
Doppelter Kehrwert ergibt die Ausgangsgröße.

\subsection*{f) $a^{3/4}$ als Wurzelausdruck}
\begin{equation}
  a^{3/4}=\sqrt[4]{a^3} .
\end{equation}

\subsection*{g) $\frac{a^5 a^{-2}}{a^3}$}
\begin{equation}
  \frac{a^5 a^{-2}}{a^3}=\frac{a^{5-2}}{a^3}=\frac{a^3}{a^3}=a^0=1 .
\end{equation}
Gedanke: alle Potenzen (auch die im Nenner, mit negativem Vorzeichen im
Exponenten) werden kombiniert: $5+(-2)-3=0$.

\subsection*{h) $\bigl(x^{2/3}\bigr)^{-3}$ ohne negativen Exponenten}
\begin{equation}
  \bigl(x^{2/3}\bigr)^{-3}=x^{2/3\cdot(-3)}=x^{-2}=\frac1{x^2} .
\end{equation}

\begin{ergebnis}
\begin{align*}
  \text{a)}\ \tfrac18, &\quad
  \text{b)}\ 4, &\quad
  \text{c)}\ \tfrac13, &\quad
  \text{d)}\ \tfrac32, &\quad
  \text{e)}\ a^5, &\quad
  \text{f)}\ \sqrt[4]{a^3},\\
  \text{g)}\ 1, &\quad
  \text{h)}\ \tfrac1{x^2}.
\end{align*}
\end{ergebnis}

\begin{kontrolle}
\textbf{Numerische Kontrolle:} a) $2^{-3}=\tfrac1{8}$; c) $27^{1/3}=3$,
Kehrwert $\tfrac13$; d) $(\tfrac49)^{1/2}=\tfrac23$, Kehrwert jedes
Schritts: $\bigl(\tfrac49\bigr)^{-1/2}=\tfrac32$. \checkmark

\textbf{Stichprobe:} h) Probe mit $x=2$:
$(2^{2/3})^{-3}=2^{-2}=\tfrac14=\tfrac1{2^2}$. \checkmark

\textbf{Grenzfall:} g) $a^5 a^{-2}/a^3=1$ für jedes $a>0$ --- die
Exponentenbilanz $0$ bedeutet, die Basis kürzt sich komplett heraus. Der
(allgemeine) Grenzfall ``$\ldots$'' der Potenzgesetze: für $a\to0$
w\"arden negative Exponenten unendlich, was man als Prozess sehen kann,
aber nicht \"uberbewerten darf.
\end{kontrolle}

\begin{merksatz}
Negative Exponenten sind Kehrwerte, gebrochene Exponenten Wurzeln; mit
den Potenzgesetzen lassen sich Brüche wie $a^5a^{-2}/a^3$ auf ein Glied
zurückführen. Die Probe per Zahlenwert (hier $x=2$) bestätigt die
Umformung.
\end{merksatz}
